\section{Overview, Architecture of the Proof, and Top-Level Reduction}
\label{sec:ch01-overview-and-spine}

This paper gives a self-contained, human-readable mathematical exposition of a
complete formal proof of the Generalized Riemann Hypothesis (GRH) for all
Dirichlet $L$-functions and, as an immediate specialization, the Riemann
Hypothesis (RH) for the Riemann zeta function $\zeta(s)$.
The underlying formal development in Lean~4 comprises $175{,}187$ source
modules totaling $15{,}531{,}759$ lines of code; exact declaration-level
dependency analysis rooted at the five headline theorems of \leanmod{Main}
identifies a reachable proof cone of $145{,}276$ modules, $799{,}064$
declarations, and $10{,}432{,}497$ declaration body lines.
Every mathematical claim presented in this manuscript is anchored to its
corresponding kernel-checked Lean constant via inline pointers of the form
\leanlink{riemannHypothesis}, allowing the reader to pass effortlessly
between classical mathematical prose and the formal source.

% -----------------------------------------------------------------------------
\subsection{Statement of the Main Results}
\label{subsec:ch01-main-results}

For a modulus $N \ge 1$ and a Dirichlet character
$\chi : (\Z/N\Z)^{\times} \to \C^{\times}$ (extended by zero to integers not
coprime to $N$), let $L(s, \chi)$ denote the analytic continuation of the
Dirichlet series $\sum_{n=1}^{\infty} \chi(n) n^{-s}$ ($\re(s) > 1$), and let
$\zeta(s) = L(s, \mathbf{1}_1)$ denote the Riemann zeta function.
The top-level module \leanmod{Statements} formulates the four canonical
zero-free and critical-line predicates
\leanlink{DirichletLNonvanishingRightOfHalf},
\leanlink{DirichletLZerosOnLine},
\leanlink{PrimitiveDirichletLZerosOnLine}, and
\leanlink{RiemannZetaZerosOnLine} directly in terms of Mathlib's definitions of
$L(s, \chi)$ and $\zeta(s)$, and \leanmod{Main} proves all five headline
theorems unconditionally.

\begin{theorem}[Right-half-plane nonvanishing for Dirichlet $L$-functions]
\label{thm:ch01-half-plane}
\leanlink{LFunction_ne_zero_of_half_lt_re}
For every integer $N \ge 1$, every Dirichlet character $\chi \pmod{N}$, and
every complex number $s \in \C$ with $\re(s) > \tfrac{1}{2}$ satisfying either
$\chi \neq 1$ or $s \neq 1$, one has
\[
  L(s, \chi) \neq 0.
\]
\end{theorem}

\begin{proof}
In the formal development, the target predicate
\leanlink{DirichletLNonvanishingRightOfHalf} is definitionally identical to
\leanlink{UpperTheory.EveryCharacter}.
The theorem is established by \leanlink{allChi} in
\leanmod{Supplements.Kernels.MainTheorems}, which instantiates
\leanlink{UpperTheory.FinalTheorems.everyCharacter} from
Section~\ref{sec:galois-cohomology-assembly}.
\end{proof}

Combining right-half-plane nonvanishing with primitive-character reduction,
nonvanishing of the archimedean gamma factor, and the functional equation of
completed Dirichlet $L$-functions confines every nontrivial zero in the
critical strip to the critical line $\re(s) = \tfrac{1}{2}$.

\begin{theorem}[Generalized Riemann Hypothesis for all Dirichlet $L$-functions]
\label{thm:ch01-grh-all}
\leanlink{LFunction_eq_zero_re}
For every integer $N \ge 1$, every Dirichlet character $\chi \pmod{N}$, and
every $s \in \C$ in the critical strip $0 < \re(s) < 1$,
\[
  L(s, \chi) = 0 \implies \re(s) = \frac{1}{2}.
\]
\end{theorem}

\begin{proof}
The predicate \leanlink{DirichletLZerosOnLine} coincides with
\leanlink{Claim.Strip}.
By \leanlink{claimAll}, applying \leanlink{Claim.strip} to the half-plane
nonvanishing result \leanlink{allChi} (Proposition~\ref{prop:ch01-claim-strip}
below) yields the conclusion for all moduli $N \ge 1$ and all characters
$\chi \pmod{N}$, whether primitive or imprimitive.
\end{proof}

\begin{corollary}[Primitive Dirichlet $L$-functions]
\label{cor:ch01-grh-primitive}
\leanlink{LFunction_eq_zero_re_of_isPrimitive}
For every integer $N \ge 1$, every primitive Dirichlet character
$\chi \pmod{N}$, and every $s \in \C$ with $0 < \re(s) < 1$ such that
$L(s, \chi) = 0$, one has $\re(s) = \tfrac{1}{2}$.
\end{corollary}

\begin{proof}
The statement \leanlink{PrimitiveDirichletLZerosOnLine} is proved directly by
\leanlink{critAll}, which invokes the direct functional-equation reflection
\leanlink{Formalism.VW.critLine} on \leanlink{allChi}.
\end{proof}

\begin{corollary}[The Riemann Hypothesis]
\label{cor:ch01-rh}
\leanlink{riemannHypothesis}
Every nontrivial zero $s \in \C$ of the Riemann zeta function $\zeta(s)$---that
is, every $s \in \C$ with $\zeta(s) = 0$, $s \neq -2(n+1)$ for all $n \in \N$,
and $s \neq 1$---lies on the critical line $\re(s) = \tfrac{1}{2}$.
\end{corollary}

\begin{proof}
This is established by \leanlink{rhOfAll}, which applies
\leanlink{Formalism.VW.rh_of_every} to \leanlink{allChi} (specializing
$\chi = \mathbf{1}_1$ and reflecting via \leanlink{Formalism.VW.rh_of_right}).
\end{proof}

\begin{remark}[Equivalent formalizations of RH]
\label{rem:ch01-rh-equiv}
The explicit zero-localization formulation \leanlink{riemannZeta_eq_zero_re}
for \leanlink{RiemannZetaZerosOnLine} is definitionally identical to
\leanlink{riemannHypothesis}, and is independently witnessed as the third
component of \leanlink{claimAll} via \leanlink{Claim.zeta}.
\end{remark}

% -----------------------------------------------------------------------------
\subsection{Conceptual Architecture and Reader's Roadmap}
\label{subsec:ch01-roadmap}

At its conceptual core, the proof realizes a number-field analogue of
Stepanov's polynomial-multiplicity method for curves over finite fields.
In characteristic $p$, high-order vanishing of an auxiliary polynomial at a
large set of rational points converts an analytic zero-counting problem into an
algebraic lower bound.
Over $\Q$, where no literal differentiation operator acts on prime numbers,
high-order vanishing is replaced by a tower of iterated definite unitary theta
lifts whose automorphic periods are simultaneously bounded from above by
Arakelov slope inequalities and positivity determinants, and bounded from below
by integral dual Weyl divisibility and Moriwaki's arithmetic bigness theorem.
Closing the resulting global inequality requires controlling the Galois and
Iwasawa-theoretic obstructions that appear in the Selmer windows attached to
cyclotomic and CM extensions.

The proof therefore interlocks four deep mathematical pillars before closing at
the top-level reduction spine.
The remainder of this manuscript is organized so that each pillar can be read
both in logical sequence and as a self-contained mathematical unit:

\begin{enumerate}
  \item \textbf{Section~\ref{sec:ch01-overview-and-spine} (Overview and
  Top-Level Reduction Spine).}
  States the five main theorems, establishes global notation and conventions,
  and traces the exact formal reduction chain from \leanmod{Main} down to the
  two interface inputs entering from the automorphic and Galois-theoretic
  layers: strong Galois local--global compatibility and the Selmer-window
  bound.

  \item \textbf{Section~\ref{sec:ch02-explicit-formula-theta-lifts} (Explicit
  Formulas, Positivity Determinants, and Iterated Theta Lifts).}
  Develops the analytic-automorphic engine that drives zero-free regions: the
  critical-line reflection mechanism, the Selmer-window reduction
  (\leanlink{UpperTheory.Top.VWin}), Euclidean and Hermitian Arakelov lattice
  slopes and determinant monotonicity, iterated definite unitary theta lifts
  and Bessel period integrals, integral dual Weyl divisibility via standard
  bideterminant bases, and unconditional arithmetic bigness of Moriwaki type.

  \item \textbf{Section~\ref{sec:ch03-rtf-jacquet-rallis-analysis}
  (Relative Trace Formulas, Jacquet--Rallis Matching, and Spectral Analysis).}
  Supplies the harmonic-analytic and representation-theoretic inputs comparing
  automorphic periods on definite unitary groups against periods on general
  linear groups.
  Its centerpiece is an unconditional all-ranks proof of the Jacquet--Rallis
  fundamental lemma via Wei Zhang's Fourier-transform method, together with
  inert principal-congruence matching, H\"ormander's weighted $L^2$
  $\bar{\partial}$-estimates on complex atlases, sharp restriction and
  Hermitian curvature bounds, $\GL_3$ cumulant estimates, and Harish-Chandra
  $(\mathfrak{g}, K)$-module spectral geometry.
  The underlying local automorphic $L$-function interface is likewise built
  from scratch in \leanmod{GroundTheory.AutomorphicForms}, including
  Godement--Jacquet local zeta integrals
  \leanlink{GroundTheory.Automorphic.GodementJacquet.zetaIntegral},
  their fractional zeta ideals
  \leanlink{GroundTheory.Automorphic.GodementJacquet.zetaIdeal},
  and local $L$-polynomials and $L$-factors
  \leanlink{GroundTheory.Automorphic.GodementJacquet.LPolynomial} and
  \leanlink{GroundTheory.Automorphic.GodementJacquet.LFunction}.

  \item \textbf{Section~\ref{sec:abelian-varieties-weil} (Abelian Varieties,
  the Theorem of the Cube, Isogenies, and Weil Bounds).}
  Constructs the algebraic- and arithmetic-geometric foundations required both
  for unitary PEL Shimura varieties and for Arakelov intersection bounds on
  their toroidal compactifications.
  This is the single largest portion of the formal dependency graph: spreading
  out and the Theorem of the Cube over arbitrary bases, multiplication-by-$n$
  isogenies, torsion cardinality and $p$-divisible groups, characteristic
  polynomials of endomorphisms and norm-degree identities, characteristic-zero
  resolution of singularities, and multiplicative height bounds of Weil and
  Mahler type.

  \item \textbf{Section~\ref{sec:galois-cohomology-assembly} (Galois
  Local--Global Compatibility, Cohomological Obstructions, and Closing the
  Proof of GRH).}
  Brings all preceding layers together in the upper arithmetic theory:
  \'etale local acyclicity over strictly Henselian bases, perfectoid tilting
  and Fontaine's period ring $A_{\mathrm{inf}}$, finite \'etale level models on
  unitary PEL Shimura varieties, strong Galois local--global compatibility for
  regular algebraic conjugate-self-dual cuspidal automorphic representations of
  $\GL_a(\A_E)$ (\leanlink{UpperTheory.Gal.Int.galoisLocalGlobalStrong_recParameter}),
  local Galois Euler characteristics, Wach dual-fibre lattices and Bost's slope
  inequality yielding the weakened Vandiver window
  (\leanlink{UpperTheory.One.vwin_of_row}), and the final assembly of
  \leanlink{UpperTheory.FinalTheorems.everyCharacter}.
\end{enumerate}

Table~\ref{tab:ch01-architecture} summarizes the distribution of the
$145{,}276$ reachable modules and $10{,}432{,}497$ reachable declaration body
lines across the major namespaces of the formal proof, indicating where each
pillar is treated in this manuscript.

\begin{table}[ht]
\centering
\caption{Quantitative architecture of the reachable formal proof cone rooted at \leanmod{Main} ($145{,}276$ reachable modules, $799{,}064$ declarations, $10{,}432{,}497$ declaration body LOC).}
\label{tab:ch01-architecture}
{\small
\begin{tabular}{@{}lrrrl@{}}
\toprule
Primary Lean package & Reachable mods & Reachable decls & Body LOC & Exposition \\
\midrule
\leanmod{Development.AbelianVarietyTheory} & 79,657 & 337,791 & 4,620,502 & Section~\ref{sec:abelian-varieties-weil} \\
\leanmod{Development.AutomorphicArithmetic} & 16,657 & 78,518 & 1,262,457 & Sections~\ref{sec:ch01-overview-and-spine}, \ref{sec:ch02-explicit-formula-theta-lifts}, \ref{sec:galois-cohomology-assembly} \\
\leanmod{Development.AutomorphicSpectrum} & 13,632 & 57,399 & 775,251 & Section~\ref{sec:ch03-rtf-jacquet-rallis-analysis} \\
\leanmod{Development.WeilBound} & 6,717 & 50,329 & 612,548 & Section~\ref{sec:abelian-varieties-weil} \\
\leanmod{Formalism.LocalToGlobal} & 4,824 & 53,640 & 560,221 & Section~\ref{sec:galois-cohomology-assembly} \\
\leanmod{GroundTheory.HeightsAndTheta} & 3,916 & 18,280 & 252,201 & Section~\ref{sec:ch02-explicit-formula-theta-lifts} \\
\leanmod{Formalism.PeriodIntegrals} & 3,360 & 24,424 & 272,263 & Section~\ref{sec:ch02-explicit-formula-theta-lifts} \\
\leanmod{Formalism.SharpRestriction} & 2,725 & 36,897 & 427,095 & Section~\ref{sec:ch03-rtf-jacquet-rallis-analysis} \\
\leanmod{ReciprocityTheory.ReductiveGroups} & 2,444 & 21,610 & 272,374 & Section~\ref{sec:ch03-rtf-jacquet-rallis-analysis} \\
\leanmod{Formalism.HermitianCurvature} & 1,778 & 29,020 & 310,170 & Section~\ref{sec:ch03-rtf-jacquet-rallis-analysis} \\
\leanmod{Development.EtaleCohomology} & 1,304 & 9,127 & 73,462 & Section~\ref{sec:galois-cohomology-assembly} \\
\leanmod{Development.ResolutionOfSingularities} & 497 & 7,590 & 116,064 & Section~\ref{sec:abelian-varieties-weil} \\
\leanmod{GroundTheory.AutomorphicForms} & 476 & 7,970 & 91,329 & Sections~\ref{sec:ch01-overview-and-spine}, \ref{sec:ch03-rtf-jacquet-rallis-analysis} \\
All remaining reachable packages & 7,289 & 66,469 & 786,560 & Sections~\ref{sec:ch02-explicit-formula-theta-lifts}--\ref{sec:galois-cohomology-assembly} \\
\midrule
\textbf{Total reachable proof cone} & \textbf{145,276} & \textbf{799,064} & \textbf{10,432,497} & Sections~1--5 \\
\bottomrule
\end{tabular}
}
\end{table}

% -----------------------------------------------------------------------------
\subsection{Notation and Conventions}
\label{subsec:ch01-notation}

Throughout the paper we adhere to standard conventions of analytic number
theory, automorphic forms, and arithmetic geometry, summarized in
Table~\ref{tab:ch01-notation}.
Whenever a formal declaration is cited, a bracketed anchor such as
\texttt{[Const]} identifies the exact kernel-checked Lean constant in the
reachable cone of \leanmod{Main}, and \leanmod{Mod.Path} identifies a Lean
module path.
Appendix~\ref{app:lean-concordance} collects a concordance mapping every
anchored theorem, proposition, lemma, corollary, and definition in the text to
its defining module and formal proof length.

\begin{table}[ht]
\centering
\caption{Global notation and conventions used across Sections~1--5.}
\label{tab:ch01-notation}
{\small
\begin{tabular}{@{}p{0.26\textwidth}p{0.68\textwidth}@{}}
\toprule
Symbol & Meaning \\
\midrule
$\N, \Z, \Q, \R, \C, \F_q$ & Nonnegative integers, integers, rationals, reals, complexes, and finite field of order $q$. \\
$\A_F, \A_F^{\times}, \mathcal{O}_F$ & Adeles, ideles, and ring of integers of a global or local field $F$. \\
$\chi \pmod{N}$, $\chi^{\circ}$, $f_{\chi}$ & Dirichlet character modulo $N \ge 1$, its associated primitive character, and its conductor $f_{\chi} \mid N$. \\
$L(s, \chi)$, $\Lambda(s, \chi)$, $\gamma_{\chi}(s)$ & Dirichlet $L$-function, completed $L$-function $\Lambda(s, \chi) = \gamma_{\chi}(s) L(s, \chi)$, and archimedean gamma factor $\Gamma_{\R}(s + a_{\chi})$. \\
$\zeta(s)$, $\Gamma_{\R}(s)$ & Riemann zeta function $\zeta(s) = L(s, \mathbf{1}_1)$ and $\Gamma_{\R}(s) = \pi^{-s/2}\Gamma(s/2)$. \\
$E/F$, $x \mapsto \bar{x}$, $\eta$, $\psi$ & Quadratic extension (CM global or unramified local), nontrivial Galois involution, associated quadratic character, and additive character. \\
$\GL_n$, $\SL_n$, $\U(V)$, $\U(n)$ & General linear, special linear, and (definite or quasi-split) unitary groups. \\
$Z(\Phi, f, s)$, $L(s, \pi)$ & Local Godement--Jacquet zeta integral and local $L$-factor of an irreducible admissible representation $\pi$ of $\GL_n(F)$. \\
$A/S$, $A^{\vee}$, $A[n]$, $T_{\ell}(A)$ & Abelian scheme over a base scheme $S$, dual abelian scheme, $n$-torsion group scheme, and $\ell$-adic Tate module. \\
$\Gal(\bar{K}/K)$, $H^i_{\mathrm{\acute{e}t}}$ & Absolute Galois group of $K$ and $\ell$-adic \'etale cohomology. \\
$A_{\mathrm{inf}}, A_{\mathrm{cris}}, B_{\mathrm{dR}}^+$ & Fontaine's $p$-adic period rings attached to a perfectoid tilt. \\
$H^1(\mathcal{O}_E[1/\ell], \Q_\ell/\Z_\ell(-e))$ & Global Galois cohomology group governing the weakened Vandiver / Selmer-window obstruction \leanlink{UpperTheory.Top.VWin}. \\
\bottomrule
\end{tabular}
}
\end{table}

% -----------------------------------------------------------------------------
\subsection{The Top-Level Formal Reduction Spine}
\label{subsec:ch01-spine}

Before entering the detailed analytic, automorphic, geometric, and
cohomological theories of Sections~2--5, we record the exact logical spine that
reduces all five headline results of \leanmod{Main} to the two upper-theory
inputs proved in Sections~2 and~5.
Reading from bottom to top, the spine splits cleanly into two halves: an
\emph{arithmetic assembly} producing right-half-plane nonvanishing for every
Dirichlet character, and an \emph{analytic reflection} deducing critical-line
localization for all Dirichlet $L$-functions and $\zeta(s)$.

\subsubsection{Arithmetic assembly of right-half-plane nonvanishing}

The single bottleneck hub through which the entire formalization passes is the
unconditional theorem \leanlink{UpperTheory.FinalTheorems.everyCharacter} in
\leanmod{Development.AutomorphicArithmetic.CharacterBoxes.MainConclusion}.
Its proof factors through two transparent steps: a conditional reduction from
strong Galois local--global compatibility, and the unconditional supply of that
compatibility from unitary theta lifts and geometric realizations on Shimura
varieties.

\begin{proposition}[Conditional half-plane nonvanishing from strong local--global compatibility]
\label{prop:ch01-every-of-row}
\leanlink{UpperTheory.FinalTheorems.everyCharacter_of_row}
Assume that strong Galois local--global compatibility holds for the local
Langlands correspondence predicate \leanlink{UpperTheory.Zero.LLC.IsRecT} on
regular algebraic conjugate-self-dual cuspidal automorphic representations
(\leanlink{UpperTheory.GalStar.Row0b}).
Then \leanlink{UpperTheory.EveryCharacter} holds: for every $N \ge 1$, every
Dirichlet character $\chi \pmod{N}$, and every $s \in \C$ with
$\re(s) > \tfrac{1}{2}$ (and $\chi \neq 1$ or $s \neq 1$), one has
$L(s, \chi) \neq 0$.
\end{proposition}

\begin{proof}
Given the strong local--global hypothesis $h$ for
\leanlink{UpperTheory.Zero.LLC.IsRecT}, the proof combines three modular
theorems:
\begin{enumerate}
  \item existence of the local Langlands parameter correspondence
  \leanlink{UpperTheory.Zero.LLC.recTExists} (built in
  \leanmod{Development.AutomorphicArithmetic.Foundations.LocalLanglands.LocalCentralUnits.RecTExists}
  from cuspidal exhaustiveness and uniqueness of Langlands extensions);
  \item the weakened Vandiver / Selmer-window obstruction bound
  \leanlink{UpperTheory.One.vwin_of_row}, which feeds $h$ together with
  definite unitary base change
  \leanlink{UpperTheory.BaseChangeIdentities.baseChangeN} into
  \leanlink{UpperTheory.One.Join.automorphic_window_of_inputs} to establish
  \leanlink{UpperTheory.Top.VWin} (proved in
  Section~\ref{sec:galois-cohomology-assembly}); and
  \item the global explicit-formula and theta-lift closure
  \leanlink{UpperTheory.LiteratureTheorems.every_character_of_literature_of_vWin},
  which combines the automorphic literature inputs, Moriwaki bigness, and
  \leanlink{UpperTheory.Top.VWin} to deduce \leanlink{UpperTheory.EveryCharacter}
  (Sections~\ref{sec:ch02-explicit-formula-theta-lifts}--\ref{sec:abelian-varieties-weil}). \qedhere
\end{enumerate}
\end{proof}

\begin{theorem}[Unconditional right-half-plane nonvanishing]
\label{thm:ch01-every-character}
\leanlink{UpperTheory.FinalTheorems.everyCharacter}
The universal right-half-plane nonvanishing statement
\leanlink{UpperTheory.EveryCharacter} holds unconditionally.
\end{theorem}

\begin{proof}
By Proposition~\ref{prop:ch01-every-of-row}
(\leanlink{UpperTheory.FinalTheorems.everyCharacter_of_row}), it suffices to
supply the strong Galois local--global compatibility theorem
\leanlink{UpperTheory.Gal.Int.galoisLocalGlobalStrong_recParameter}.
In \leanmod{Development.AutomorphicArithmetic.Foundations.GaloisRepresentations.IntegralStructures.RowSrc},
this theorem is obtained by invoking
\leanlink{UpperTheory.Gal.Int.galoisRepresentations_of_thetaLift} on the
theta-lift kernel package \leanlink{UpperTheory.TP.K.tLiftK}, which internally
pairs the weight-one geometric realization
\leanlink{UpperTheory.Gal.Int.Geo.geoOne_holds} with the theta-parseval lift
to construct the required integral Galois representations
(Section~\ref{sec:galois-cohomology-assembly}).
\end{proof}

\subsubsection{Analytic reflection to the critical line}

Once \leanlink{UpperTheory.EveryCharacter} is established (re-exported as
\leanlink{allChi} in \leanmod{Supplements.Kernels.MainTheorems}), the passage
from $\re(s) > \tfrac{1}{2}$ to the full critical strip $0 < \re(s) < 1$ is
elementary and completely self-contained.
The formalization implements this reflection cleanly in
\leanmod{Development.AutomorphicArithmetic.CharacterBoxes.ClaimStrip} (covering
all characters, primitive or imprimitive, as well as $\zeta(s)$) and in
\leanmod{Formalism.GammaFactorNonvanishing.CritLine} (giving the direct path
for primitive characters and $\zeta(s)$, detailed further in
Subsection~\ref{subsec:ch02-critical-line}).

\begin{lemma}[Euler-factor nonvanishing and reduction to primitive characters]
\label{lem:ch01-prim-zero}
\leanlink{Claim.prim_zero}
Let $N \ge 1$ and $\chi \pmod{N}$ be a Dirichlet character of conductor
$f_{\chi} \ge 1$ with associated primitive character $\chi^{\circ} \pmod{f_{\chi}}$.
If $s \in \C$ satisfies $0 < \re(s) < 1$ and $L(s, \chi) = 0$, then
$L(s, \chi^{\circ}) = 0$.
\end{lemma}

\begin{proof}
Since $\re(s) < 1$, we have $s \neq 1$, so the level-change identity relates
$L(s, \chi)$ to $L(s, \chi^{\circ})$ via a finite product of local Euler
factors over primes $p \mid N$ not dividing $f_{\chi}$:
\[
  L(s, \chi) = L(s, \chi^{\circ}) \prod_{p \mid N} \bigl(1 - \chi^{\circ}(p) p^{-s}\bigr).
\]
For every prime $p$ and every $s \in \C$ with $\re(s) > 0$, we have
$|\chi^{\circ}(p)| \le 1$ and $|p^{-s}| = p^{-\re(s)} < 1$, hence
$|\chi^{\circ}(p) p^{-s}| < 1$ and therefore $1 - \chi^{\circ}(p) p^{-s} \neq 0$
(\leanlink{Claim.factor_ne_zero}).
Since none of the finite Euler factors can vanish, $L(s, \chi) = 0$ forces
$L(s, \chi^{\circ}) = 0$.
\end{proof}

\begin{lemma}[Functional-equation reflection for primitive characters]
\label{lem:ch01-prim-reflect}
\leanlink{Claim.prim_reflect}
Let $M \ge 1$ and let $\psi \pmod{M}$ be a primitive Dirichlet character.
If $s \in \C$ satisfies $0 < \re(s) < 1$ and $L(s, \psi) = 0$, then
$L(1 - s, \psi^{-1}) = 0$.
\end{lemma}

\begin{proof}
Because $0 < \re(s) < 1$, both $s$ and $1 - s$ are nonzero and have strictly
positive real parts.
By parity of $\psi$, the archimedean gamma factor $\gamma_{\psi}(s)$ equals
either $\Gamma_{\R}(s)$ or $\Gamma_{\R}(s + 1)$, neither of which vanishes in
the right half-plane $\re(s) > 0$ (\leanlink{Claim.gamma_ne_zero}).
From $L(s, \psi) = \Lambda(s, \psi)/\gamma_{\psi}(s) = 0$ we obtain
$\Lambda(s, \psi) = 0$.
Since $\psi^{-1}$ is again primitive of conductor $M$, the functional equation
of completed Dirichlet $L$-functions gives $\Lambda(1 - s, \psi^{-1}) = 0$, and
dividing by $\gamma_{\psi^{-1}}(1 - s)$ yields $L(1 - s, \psi^{-1}) = 0$.
\end{proof}

\begin{proposition}[From right-half-plane nonvanishing to the critical line]
\label{prop:ch01-claim-strip}
\leanlink{Claim.strip}
Assume \leanlink{Claim.HalfPlane} (that is, $L(s, \chi) \neq 0$ whenever
$\re(s) > \tfrac{1}{2}$ and either $\chi \neq 1$ or $s \neq 1$).
Then \leanlink{Claim.Strip} holds: for every $N \ge 1$, every Dirichlet
character $\chi \pmod{N}$, and every $s \in \C$ with $0 < \re(s) < 1$ and
$L(s, \chi) = 0$, one has $\re(s) = \tfrac{1}{2}$.
\end{proposition}

\begin{proof}
Let $0 < \re(s) < 1$ and $L(s, \chi) = 0$; in particular $s \neq 0$ and
$s \neq 1$.
By trichotomy on $\re(s)$ relative to $\tfrac{1}{2}$:
\begin{itemize}
  \item If $\re(s) > \tfrac{1}{2}$, then $L(s, \chi) = 0$ immediately
  contradicts \leanlink{Claim.HalfPlane} (since $s \neq 1$).
  \item If $\re(s) < \tfrac{1}{2}$, then Lemma~\ref{lem:ch01-prim-zero}
  (\leanlink{Claim.prim_zero}) gives $L(s, \chi^{\circ}) = 0$, and
  Lemma~\ref{lem:ch01-prim-reflect} (\leanlink{Claim.prim_reflect}) reflects
  this to $L(1 - s, (\chi^{\circ})^{-1}) = 0$.
  Since $\re(1 - s) = 1 - \re(s) > \tfrac{1}{2}$ and $1 - s \neq 1$ (because
  $s \neq 0$), this contradicts \leanlink{Claim.HalfPlane}.
\end{itemize}
Hence $\re(s) = \tfrac{1}{2}$.
\end{proof}

\begin{proposition}[Specialization to the Riemann zeta function]
\label{prop:ch01-claim-zeta}
\leanlink{Claim.zeta}
Assume \leanlink{Claim.HalfPlane}.
Then \leanlink{Claim.Zeta} holds: every zero $s \in \C$ of $\zeta(s)$ with
$s \neq -2(n+1)$ ($n \in \N$) and $s \neq 1$ satisfies $\re(s) = \tfrac{1}{2}$.
\end{proposition}

\begin{proof}
Suppose $\zeta(s) = 0$, $s \neq 1$, and $s$ is not a negative even integer.
Since $\zeta(0) = -\tfrac{1}{2} \neq 0$, we have $s \neq 0$.
If $\re(s) \ge 1$, this contradicts the classical nonvanishing of $\zeta$ on
$\re(s) \ge 1$.
If $0 < \re(s) < 1$, then $L(s, \mathbf{1}_1) = \zeta(s) = 0$, so
Proposition~\ref{prop:ch01-claim-strip} (\leanlink{Claim.strip}) gives
$\re(s) = \tfrac{1}{2}$.
Finally, if $\re(s) \le 0$, then $s$ is neither $0$ nor a negative even
integer, so $\Gamma_{\R}(s) \neq 0$; hence the completed zeta function
$\Lambda(s)$ vanishes, and the functional equation $\Lambda(1 - s) = \Lambda(s)$
gives $\zeta(1 - s) = 0$ with $\re(1 - s) \ge 1$, again contradicting
nonvanishing on the closed half-plane $\re \ge 1$.
\end{proof}

The top-level kernel module \leanmod{Supplements.Kernels.MainTheorems} packages
Theorem~\ref{thm:ch01-every-character} as \leanlink{allChi}, combines it with
Propositions~\ref{prop:ch01-claim-strip} and~\ref{prop:ch01-claim-zeta} into
the triple \leanlink{claimAll}, and invokes \leanlink{Formalism.VW.critLine}
and \leanlink{Formalism.VW.rh_of_every} to produce \leanlink{critAll} and
\leanlink{rhOfAll}.
Finally, \leanmod{Main} exports the five headline results
\leanlink{riemannHypothesis}, \leanlink{riemannZeta_eq_zero_re},
\leanlink{LFunction_ne_zero_of_half_lt_re}, \leanlink{LFunction_eq_zero_re},
and \leanlink{LFunction_eq_zero_re_of_isPrimitive}.
Thus the entire proof of GRH and RH is reduced with zero logical gaps to
establishing the analytic-automorphic, geometric, and Galois-theoretic engines
exposed in Sections~\ref{sec:ch02-explicit-formula-theta-lifts}
through~\ref{sec:galois-cohomology-assembly}.
